\section{Abstract}

L'obiettivo del seguente progetto è la realizzazione di una base di dati per la gestione di
un concessionario di automobili del gruppo BMW e la relativa officina meccanica.
In particolare, si focalizza sulla gestione delle dinamiche di vendita e assistenza dei
veicoli, sulla gestione dei clienti e del personale e sul tracciamento
dei pezzi di ricambio necessari agli interventi.

L'architettura del database rispecchia le dinamiche del mercato automobilistico moderno.
Il sistema mappa il catalogo dei modelli commerciali e le relative specifiche tecniche,
separando le caratteristiche costruttive dei propulsori dalle configurazioni software di 
potenza e dalla gestione degli accessori opzionali. Viene modellato il parco auto
disponibile, ponendo una separazione tra veicoli nuovi, con l'elenco degli accessori 
opzionali presenti, e usati, con lo storico dei proprietari e degli interventi di
manutenzione effettuati.

Una sezione cruciale del sistema è dedicata all'area assistenza, in cui vengono tracciati 
i flussi operativi dell'officina: dalla registrazione delle commesse di assistenza,
all'allocazione del personale tecnico specializzato, fino al monitoraggio
dei consumi di magazzino.

Il database centralizza infine la gestione della clientela e il tracciamento delle
transazioni economiche finalizzate dall'organico commerciale.

La base di dati è progettata per ottimizzare il flusso di lavoro all'interno
dell'officina, garantendo una gestione efficiente delle risorse e una migliore esperienza 
per i clienti.